\section{Eigenvalues of a finite motif}
\label{sec:h-levels}

A motif on \(n\) vertices has at most \(n\) eigenvalues of \(H=-A\).
Those below the essential-spectrum edge are stable against the continuum
\eqref{eq:dos}.
Those inside the band are resonances whose widths track \(\rho(E)\).
Isolated \(K_n\) has one level \(-(n-1)\) and an \((n-1)\)-fold level \(+1\).
Only the ground level can lie below \(-8\), and only for \(n\ge 10\).
The level \(+1\) lies in \([-8,8]\).
Coupling to the lattice can split that multiplet into resonances.
It pushes a vector below the edge only if some trial vector says so;
the all-ones direction is the one that scores \(-(n-1)\).
Commutant of the motif automorphism group with \(H\) gives the selection rules
among those eigenvalues.
This spectrum is the internal structure whose ground gap is
Section~\ref{sec:h-mass}.
The curvature \(m^*=1/4\) of the continuum above the edge is not one of
these eigenvalues (Section~\ref{sec:masses}).
On an eigenstate the level also fixes the bond sum at each vertex,
Lemma~\ref{lem:starsum}.
The flat band forces that sum to vanish and does not force each bond to
vanish.
